\documentclass[twoside,12pt]{dscproblemset}

\usepackage[margin=1in]{geometry}
\usepackage{amsmath}
\usepackage{amssymb,amsthm}
\usepackage{fancyhdr}
\usepackage{nicefrac}
\usepackage{tikz}
\usepackage{pgfplots}
\usepackage[outputdir=build]{minted}
\usetikzlibrary{quotes,angles,positioning,arrows.meta}
\usetikzlibrary{calc}

% configuration
% ------------------------------------------------------------------------------

% page headers only on odd pages
\pagestyle{fancy}
\fancyhead{}
\fancyhead[RO]{PID or Name: \rule{3in}{.5pt}}
\renewcommand{\headrulewidth}{0pt}

% commands
% ------------------------------------------------------------------------------

\newcommand{\python}[1]{\mintinline{python}{#1}}

% ------------------------------------------------------------------------------


\begin{document}

\examfrontpage{
    % course
    DSC 40B
}{
    % exam name
    Redemption Midterm 01
}{
    % exam date
    March 21, 2023
}{
    % instructions
    \textbf{Instructions:}
    \begin{itemize}
        \item Write your solutions to the following problems in the spaces provided.
        \item No calculators are permitted, but a few pages of notes are.
        \item Write your name or PID at the top of each sheet in the space provided.
        \item If asked to state a time complexity, do so using asymptotic notation in as simplest terms possible.
    \end{itemize}

    \begin{center} 
        {
        \fontsize{20pt}{30pt}\selectfont
        Redemption Midterm\\[1em]
        \fontsize{80pt}{30pt}\selectfont
        01
        }
    \end{center}
    \vspace{2em}

}

\begin{probset}

    % Time Complexity of Loops
    % ======================================================================================

    % independent loops
    \inputproblem{01}

    % dependent while-loops, arithmetic
    \inputproblem{02}

    % ugly bounds
    \inputproblem{03}

    \newpage

    % built-in functions
    \inputproblem{04}

    % best case
    \inputproblem{05}

    % worst case
    \inputproblem{06}

    \newpage

    % count number of elements that are within 10 of the maximum, unsorted
    % what is a theoretical lower bound for this problem?
    \inputproblem{07}

    % expected time
    \inputproblem{09}

    \newpage

    % expected time
    \inputproblem{08}

    % state recurrence
    \inputproblem{10}

    \newpage

    % solve recurrence
    \inputproblem{11}

    \newpage

    % asymptotic notation
    \inputproblem{12}

    % asymptotic notation
    \inputproblem{13}

    % asymptotic notation
    \inputproblem{14}

    % asymptotic notation
    \inputproblem{15}

    % modified binary search
    \inputproblem{16}

    \inputproblem{17}

    \newpage

    \inputproblem{18}

    \inputproblem{19}

    \inputproblem{20}

\end{probset}

\vfill
\begin{center}
    \textbf{Before turning in your exam, please check that your name is on every page.}
\end{center}


\ifshowsoln
\else
    \newpage
    \scratchpage{}
    \newpage
    \scratchpagedetach{}
    \newpage
    \scratchpagedetach{}
\fi

\end{document}
